\section{Auditoría de fuentes: el hilo conductor del Transformer es la Consolidation}

Esta clase no es un libro de texto estándar sobre el Transformer, ni una explicación página por página del artículo «Attention Is All You Need». El coautor Ashish Vaswani parte de la Dartmouth proposal de 1955 y rastrea cómo la inteligencia artificial pasó de un sueño de unified-system a fragmentarse en innumerables pipelines especializados, y cómo, gracias a distributed representation, sequence-to-sequence, attention y el aprendizaje a gran escala, volvió a encaminarse hacia la consolidation. En la segunda mitad, la perspectiva se amplía del algoritmo a position representation, long context, memory bandwidth, tools, data center y el ciclo de retroalimentación del producto.

\subsection{Límites oficiales de la clase y recuperación visual}

Esta reescritura toma como fuente del orden de la clase la grabación oficial de Stanford Online \texttt{1GbDTTK3aR4}; la clase se impartió el 7 de noviembre de 2023. La grabación incluye subtítulos \texttt{en-US} hechos por personas, de los que se analizaron 1,866 cues válidos. En la página oficial, en la descripción de la grabación y en la búsqueda por exact-title no se encontró ningún standalone deck público, por lo que la grabación en 1080p se usa como fuente visual; tras una extracción de alta exhaustividad cada tres segundos, se conservaron manualmente 43 estados didácticos independientes de 87 candidatos, y se omitieron deliberadamente los fotogramas repetidos de speaker-window, las revisitas repetidas, los bumper puros y las páginas de Thanks que aparecen varias veces.

\begin{longtable}{@{}L{0.18\textwidth}L{0.31\textwidth}L{0.40\textwidth}@{}}
\toprule
Capa de evidencia & Material local & Función en la nota \\
\midrule
Grabación oficial & 1:20:38, 43 estados didácticos independientes & Determina el orden de la clase, la explicación de los gestos, las demos, las Q\&A y la narrativa histórica \\
Subtítulos hechos por personas & 1,866 cues, transcripciones timed/clean/chunk & Recuperan la motivación, las rutas fallidas, los juicios de ingeniería, las reservas y las líneas de investigación \\
Primary sources & Seq2Seq, NMT Attention, Transformer, Music Transformer, Routing/Reformer, MQA/GQA, FlashAttention, entre otras & Verifican fórmulas, términos, objetos de los experimentos y límites temporales \\
Auditoría de cobertura y visual & manifest, selection, coverage, PDF QA & Demuestran que las imágenes y los spoken nodes están planificados, tienen una ubicación y que el producto final es legible \\
\bottomrule
\end{longtable}

\begin{warningbox}{Límites históricos y límites de la evidencia}
Este texto reconstruye la clase del 2023-11-07. Los juicios del ponente sobre future architecture, agents, tools, adaptive thinking y el ciclo cerrado del producto pertenecen a la agenda de investigación de ese momento y no equivalen a algo ya validado por productos posteriores; los primeros resultados en MT/parsing demuestran que el Transformer es eficaz en las tareas dadas, pero por sí solos tampoco demuestran que sea la arquitectura general definitiva.
\end{warningbox}

\subsection{Dartmouth proposal: el objetivo era acertado, el juicio sobre el camino fue muy erróneo}

La clase usa la Dartmouth proposal para establecer un contraste deliberado. La propuesta de 1955 ya planteaba el ambicioso objetivo de que «en principio, el aprendizaje y la inteligencia pueden describirse con precisión y una máquina puede simularlos», y además consideraba que un grupo de científicos podría lograr avances significativos en un solo verano. El problema no estaba en la ambition, sino en dos juicios: sobrestimaron la capacidad humana para escribir a mano las reglas de la inteligencia y subestimaron la escala que alcanzarían la capacidad de las máquinas, los datos y los sistemas de optimización.

\lecturefigure{slide-01-dartmouth-1955.jpg}{La Dartmouth proposal de 1955 y la visión temprana de una inteligencia unificada}{Diapositiva V001 recuperada de la grabación oficial; 00:01:27}

\begin{knowledgebox}{Lectura de la figura: primero las líneas rojas, después las personas}
Los tres puntos destacados de la proposal, a la derecha, son «la inteligencia puede describirse con precisión», «en un verano pueden lograrse avances significativos» y «el obstáculo principal no es la machine capacity». Primero compare la tercera frase con la infraestructura actual de los grandes modelos; después, las hand-written rules implícitas en la primera. La figura respalda que los objetivos de la investigación temprana tienen continuidad, pero también muestra que el camino de implementación pasó de la symbolic programming a la data-driven optimization.
\end{knowledgebox}

\teachervoice{El ponente bromea con que los investigadores de Dartmouth necesitaban un «verano muy, muy largo», y señala que el Transformer actual ya no es una máquina, sino un data center. Su juicio más serio es que todavía no hemos aprendido a escribir todas las reglas de la inteligencia, pero sí hemos aprendido a hacer que los modelos ajusten una gran cantidad de reglas a partir de los datos.}

\begin{importantbox}{La Consolidation Thesis de esta clase}
La importancia histórica del Transformer no se reduce a ser un sequence model más rápido: las funciones que antes estaban dispersas entre feature engineering, alignment, language modeling, task-specific architecture y el posprocesamiento entran cada vez más en un sistema compartido y entrenable. La Consolidation reduce el número de interfaces, pero no elimina los problemas de datos, optimización, sistemas y evaluación.
\end{importantbox}

\subsection{Resumen de la sección}

Esta sección fijó los límites de la clase y de la evidencia, y volvió a situar el Transformer dentro de la historia más larga de la IA: el objetivo de una inteligencia unificada nunca desapareció; lo que cambió fue la forma de alcanzarlo. Lo que sigue explicará la arquitectura siguiendo el hilo de «cómo los pipelines especializados fueron absorbidos por representaciones generales y grafos de cómputo entrenables», en lugar de tratar la attention como una fórmula aislada.
